\section{Periods on the original links}\label{sec:periods}

Let $B=\Def(X)$ and fix a representative of its proper flat family near
the given compact fibre. The morphism is Gorenstein after shrinking,
and its relative dualizing line commutes with base change. A generator
of $\omega_X$ extends to a generator $\Omega$ of this line, including
over the nonreduced base. Indeed $H^1(X,\omega_X)=0$ by duality and
$H^2(X,\OO_X)=0$. In a finite free complex computing the cohomology of
the relative dualizing line and its base changes, split the invertible
entries at the origin. The remaining complex has a rank-one term in
degree zero and zero terms in degrees $-1$ and $1$. Thus its generator
extends and properness permits us to remove the image of its zero
locus. Applied to $\OO$, the same argument shows that any other
generator differs by a unit on the base. Compare \cite[Lemma~3.1]{Friedman}.

On the punctured germ at $p$ a holomorphic relative three-form
determines a class on the oriented five-dimensional link. We use its
Poincar\'e dual in $V_p$. We give the details of its variation because
the possible Artin families include maps into a union of local
components which factor through none of them.

\begin{lemma}\label{lem:linkperiods}
Let $\mathcal U\to D_p$ be a marked semiuniversal family of one of the
local germs and let $D$ be a smooth parameter germ. Choose a relative
dualizing generator on $\mathcal U$ and multiply it by any holomorphic
unit on $\mathcal U\times D$. Its marked link class gives an analytic
isomorphism over $D$,
\[
       (D_p\times D,0)\longrightarrow (G_p\times D,0).
\]
It induces the cohomological link class on every Artin pullback. It
identifies the two components at a degree-six vertex and all the marked
local discriminant hyperplanes, with their ideals.
\end{lemma}
\begin{proof}
At an ordinary double point use the equation
$\sum_{j=1}^4 x_j^2=\tau$. The boundary three-cycle represents the
vanishing sphere in a smooth fibre. Its residue period is $\tau$ times
a unit. Scaling by $\sqrt{\tau}$ gives the spherical average of the
multiplier; odd terms integrate to zero and the constant term is
nonzero. This proves the assertion also in the presence of the extra
parameters.

At a degree-six vertex the analytic base has ring
$\C\{s,u,v\}/(su,sv)$ \cite[\S\S5--8]{Alt}. Work first on either
of its two smooth branches. The simultaneous modifications of
\cite[\S3, Proposition~3.8]{W5} have two marked nodes on the central
fibre over the line and three over the plane. Their total spaces are
respectively the total spaces of $\OO(-1,-1)^{\oplus2}$ over
$\PP^1\times\PP^1$ and $\OO(-1)^{\oplus3}$ over $\PP^2$. The map is
an isomorphism off the exceptional zero section, of codimension at
least two in each smooth total space. Since both maps to the branch
are flat Gorenstein, Hartogs extension for the pullback relative
generator and its inverse gives a relative generator on the incidence
space, also after taking its product with $D$. Let $q_j$ be the signed
nodal periods there. On the line both equation parameters have nonzero
derivative. On the plane the three are cyclic differences of three
diagonal entries modulo translation, so they have joint rank two at
the origin. Each $q_j$ is its equation parameter times a unit.

To identify the periods on the original link, use the cobordism between
that link and the links of the incidence nodes. On the central fibre
it is the region in the resolved neighborhood of $E_p$ obtained by
removing the inner small-resolution neighborhoods and truncating at
the outer boundary. It transports smoothly in the incidence family
by local triviality of its compact exterior with boundary. Each
$D_0\in H^2(E_p;\C)$ extends to the local resolution by retraction.
On this cobordism the boundary of the relative Poincar\'e dual of its
restriction has outer term $\operatorname{PD}(D_0|_{L_p})$ and inner
terms $(D_0\cdot C_j)\gamma_j$, where $C_j$ is the marked exceptional
curve at the inner node and $\gamma_j$ its signed link cycle. This
follows by restriction of Poincar\'e duality to each boundary component.
Stokes' theorem for the closed relative volume form gives the outer
period as $\sum_j(D_0\cdot C_j)q_j$. The circle-bundle sequence
identifies the outer class with $\sum_jq_j[C_j]$ in $H_2(E_p;\C)$.
Taking $D_0=K_{E_p}$ gives $\sum_j q_j=0$. The two or three marked
curves are the opposite pair or the alternating triple in the
introduction \cite[Proposition~3.7]{W5}. Their weighted sums identify
the sum-zero coefficient spaces with $V_{p,L}$ and $V_{p,P}$.

The resulting maps have invertible derivative along the respective
branches and vanish at their origins for all parameters in $D$.
Their inverses over $D$ therefore agree along $\{0\}\times D$.
The analytic ring of the union is the fibre product of the two branch
rings over $\OO_D$, so both maps and their inverses glue. A nonzero
plane parameter with one zero $q_j$ leaves a single ordinary double
point. The origin of either branch retains the original cone.

At a degree-five vertex the base is smooth of dimension four.
Write $E_p=\operatorname{Bl}_4\PP^2$ with line and exceptional classes
$H,e_1,\dots,e_4$. In the marking by the conic pencils put
\[
 f_i=H-e_i\ (1\leq i\leq4),\qquad f_5=2H-\textstyle\sum e_i,\qquad
 \lambda_i(c)=c\cdot f_i .
\]
These functionals identify $V_p=K_{E_p}^{\perp}$ with
$\{\sum\lambda_i=0\}\subset\C^5$. The local Grassmannian
family has five discriminant hyperplanes $\alpha_i(b)=0$ and one
residual node per vanishing $\alpha_i$ at a noncentral parameter
\cite[Proposition~3.1 and Lemma~3.2]{W5}. Its five marked boundary
cycles evaluate the link class by nonzero multiples of the $\lambda_i$.
For the homogeneous generator their periods are nonzero multiples of
$\alpha_i$ \cite[Lemma~3.4]{W5}. For a multiplier depending on $D$,
the period still vanishes on $\{\alpha_i=0\}\times D$ by the nodal
calculation on a dense open subset of that hyperplane. Holomorphic
division gives an $\alpha_i$ factor. Scaling the homogeneous generator
along a general line gives first derivative equal to the original
derivative multiplied by the value of the unit at the vertex. This is
the calculation of \cite[proof of Lemma~3.4]{W5}, with $D$ in place
of a disc. Consequently each remaining factor is a unit, the
derivative has rank four, and all five hyperplanes have the asserted
ideals.

We finish by checking the assertion over Artin rings. On the smooth
exterior of the central singularity in an Artin family over $A$, the
relative holomorphic de Rham complex resolves the constant sheaf $A$.
Indeed local relative polydisc coordinates exist by smoothness and
the radial Poincar\'e homotopy is coefficientwise. A relative closed
form thus defines a class in $H^3(L_p;\C)\otimes A$ on a small
punctured neighborhood, compatibly with restriction and base change.
For a smooth local base the analytic period just computed has this
class as its Taylor expansion. One way to see the compatibility is
to restrict the relative de Rham resolution of the inverse-image
sheaf of base functions to the compact boundary of a Milnor tube in
the smooth locus. That boundary is a differentiable product over a
small ball. Its degree-three direct image is the constant marked
cohomology tensored with the holomorphic base functions, computed
on a finite good cover of the link. Restriction of sheaves to the
boundary is exact, and in product charts the local relative de Rham
resolution commutes with evaluation in $A$. This proves the claim for
the Taylor jets, with any additional smooth parameters.

At the degree-six vertex apply this comparison to the Artin quotients
by powers of the maximal ideal of $\OO_{D_p\times D,0}$. Each such
quotient injects into the direct sum of the two corresponding branch
quotients: in the coordinates $(s,u,v)$ and the parameters of $D$ a
surviving monomial involves either $s$ or $u,v$, or neither. The
cohomological class and the Taylor jet of the glued analytic functions
agree in each branch quotient. Injectivity gives their equality on
the whole quotient. Every Artin point factors through one of these
quotients, proving the claim also on points outside either branch
functor.
\end{proof}

\begin{proposition}\label{prop:period}
The periods of $\Omega$ define an analytic morphism $P:B\to K\times_VG$.
Its differential has image $K$ and kernel $H^1(X,\Theta_X)$.
In each singular neighborhood it is the pullback of the local
isomorphism of Lemma~\ref{lem:linkperiods} over a smooth ambient germ
containing $B$.
\end{proposition}
\begin{proof}
Embed $B$ in a smooth germ $D$. Fix classifying maps $\phi_p:B\to
D_p$ with isomorphisms of the pulled-back local families. In the
family $\mathcal U_p\times D$ the restriction of the global family near $p$ is the closed
subgerm cut out by the graph of $\phi_p$. The ratio of the global
form to a pullback local generator is a holomorphic unit there.
It lifts to a holomorphic unit near the central point of
$\mathcal U_p\times D$. The preceding lemma gives a period morphism
over $D$ whose pullback is $P_p$. This construction respects the
fixed markings and records the dependence of the multiplier on global
parameters, even on a nonreduced base.

For every Artin point of $B$ the restricted global form gives one class
on the regular locus with coefficients in its base algebra. Its
restrictions to all the links are the classes just constructed.
Take the exterior in $Y$ of the local resolution neighborhoods.
For any degree-two cohomology class on $Y$, the Poincar\'e duals
of its restrictions to the boundaries together bound a relative
four-cycle in the exterior. Thus the corresponding sum of evaluations
of the global three-class is zero. With the oriented link identifications
this says precisely that the period tuple lies in the kernel
\eqref{eq:kernel}, with coefficients in the Artin algebra.
Apply this to every maximal-ideal jet of $B$. Separatedness of its
analytic local ring gives the linear equations of $K$ on $B$.
The equations of $G$ already hold on the individual local bases.

The derivative in directions tangent to $D$ at the vertex vanishes, as
all the central link classes are zero for every multiplier. The local
period derivatives on the local tangent spaces are isomorphisms.
Hence the full period map has precisely the kernel of tangent
localization. The local-to-global sequence and the marked
identification in \cite[Theorem~2.5]{W5} give that kernel and show
that the rank is $\dim K$. The linear equations just proved identify
its image with $K$ itself.
\end{proof}
